% Abstract recipe:
% 1. Set up the scientific problem: ARC as an RL domain.
% 2. State the gap: CPU-bound environments bottleneck research; no multi-algorithm baselines exist.
% 3. State the contribution: JaxARC + throughput + baseline benchmark.
% 4. State the main empirical findings: algorithm ordering, task-family variation, observation analysis.
% 5. State the implication: practical substrate for future ARC RL research.
%
% Keep to one paragraph. Avoid: solve, breakthrough, state of the art, robust transfer.

The Abstraction and Reasoning Corpus (ARC) poses a distinctive challenge for reinforcement learning: agents must learn to edit grids through sequential actions, guided only by sparse rewards, across hundreds of structurally diverse tasks. Studying RL on ARC at scale has been impractical because existing environments run on the CPU and cap throughput at roughly 60,000 steps per second. We introduce JaxARC, an open-source, JAX-native ARC environment that runs entirely on hardware accelerators, reaching over 100 million steps per second on a single H100 GPU---a $5{,}400\times$ speedup over the prior ARCLE environment at matched batch sizes. To demonstrate what this throughput enables, we provide three contributions beyond the environment itself. First, we profile the computational overhead of different observation and action wrappers, showing that answer-visible configurations substantially improve learnability while richer observation channels carry measurable throughput costs. Second, we present the first multi-algorithm RL baseline benchmark for ARC, evaluating PPO, DQN, PQN, and REINFORCE across six ARC-AGI difficulty tiers and six ConceptARC concept groups, with five seeds and 10 million steps per condition. The results reveal a consistent algorithm hierarchy---PPO dominates, followed by REINFORCE, with DQN and PQN trailing---and large performance variation across task families, from near-perfect returns on ConceptARC Count to negligible progress on ARC-AGI Hard. JaxARC, together with these baselines and analyses, provides a practical foundation for systematic RL research on abstract reasoning.